%%%%%%%%%%%%%%%%%%%%%%%%%%%%%%%%%%%%%%%%%%%%%%%%%%%%%%%%%%%%%%%%%%%%%%%%%%%%%%%%%%%%%%%%%%%%%%%%%%%%%%
%
%   Filename    : chapter_4.tex
%
%   Description : Research Methodology. The preliminary study and the calendar moved to
%                 chapter_5.tex on 14 Nov 2026 (adviser's request).
%                 Facts and numbers come from docs/methodology-and-results.md, the extraction and
%                 training code in src/, and dlsu-manuscript/notes/claim-audit.md (approved wording).
%
%%%%%%%%%%%%%%%%%%%%%%%%%%%%%%%%%%%%%%%%%%%%%%%%%%%%%%%%%%%%%%%%%%%%%%%%%%%%%%%%%%%%%%%%%%%%%%%%%%%%%%

\chapter{Research Methodology}
\label{sec:methodology}

This chapter describes how the study will be carried out. Section~\ref{sec:design} states the research design, and Sections~\ref{sec:data} to \ref{sec:evaluation} describe the data, feature extraction, frame selection rules, classifiers and frame combinations, and evaluation procedure. Section~\ref{sec:ethics} addresses ethical considerations.


%%
%% --- 4.1 Research Design --- %%
%%

\section{Research Design}
\label{sec:design}

The design answers the two research questions of Section~\ref{sec:overview}: how selecting the frames of a video clip by their content compares with selecting them by their position in the clip and with a random choice of the same number of frames, and how much representing the order and change of the selected frames adds to their average.

The preliminary study (Section~\ref{sec:prelim}) compared one position-based rule, Start-Middle-End (SME) sampling at three frames, with content-based selection at three, five, and seven frames. Because the position-based rule existed only at three frames, the effect of how frames are chosen could not be separated from the effect of how many frames are kept. The design also had no condition showing how an arbitrary choice of frames would perform.

The proposed study therefore uses a full factorial design. Four selection rules (positional, random, $k$-Changepoint, and $k$-BOCPD) are crossed with three frame budgets ($k = 3, 5, 7$) and four classifiers: a Long Short-Term Memory network (LSTM), a Multi-Layer Perceptron (MLP), a Support Vector Machine (SVM), and a Random Forest. The $k$-Changepoint rule is the content-based rule of the preliminary study, and $k$-BOCPD is added as a second definition of content (Section~\ref{sec:selection}). SME at $k = 3$ is kept as a thirteenth selection condition to connect the results to the preliminary study. Table~\ref{tab:conditions} lists the selection conditions. Within every selection condition, the selected frames are combined in four ways for the MLP, SVM, and Random Forest, and in two ways for the LSTM (Section~\ref{sec:training}). The dataset, data split, feature set, classification task, and classifier settings are the same in every condition, so conditions differ only in which frames reach the classifier and how they are combined. Figure~\ref{fig:pipeline} shows the pipeline.

\begin{table}[t]
   \centering
   \caption{Frame selection conditions. Each condition is evaluated with all four classifiers.}
   \label{tab:conditions}
   \begin{tabular}{lL{4.6cm}cL{3.4cm}}
      \toprule
      Rule & Selects frames by & $k$ & Status \\
      \midrule
      SME & First, middle, and last valid frame & 3 & Completed; to be re-extracted \\
      Positional & Center of each of $k$ equal segments & 3, 5, 7 & Proposed \\
      Random & Uniform random draw, five draws & 3, 5, 7 & Proposed \\
      $k$-Changepoint & Top-$k$ facial-landmark velocity peaks & 3, 5, 7 & Completed; to be re-extracted \\
      $k$-BOCPD & Top-$k$ beginnings of lasting changes in facial-landmark position & 3, 5, 7 & Proposed \\
      \bottomrule
   \end{tabular}
\end{table}

\begin{figure}[t]
   \centering
   \begin{tikzpicture}[
      node distance=0.45cm,
      box/.style={draw, rounded corners=2pt, align=center, text width=9.5cm, inner sep=5pt, font=\small},
      half/.style={draw, rounded corners=2pt, align=center, text width=4.3cm, minimum height=1.5cm, inner sep=5pt, font=\small},
      arr/.style={-{Stealth[length=2.2mm]}, thick}
   ]
      \node[box] (clip) {EngageNet clip (10 s) and its engagement label};
      \node[box, below=of clip] (scan) {\textbf{First pass.} Every second frame: face and pose detection, validity check, pose values, facial-landmark positions and velocity};
      \node[box, below=of scan] (select) {\textbf{Frame selection} from the valid frames: SME ($k = 3$), positional, random, $k$-Changepoint, or $k$-BOCPD, with $k \in \{3, 5, 7\}$};
      \node[box, below=of select] (extract) {\textbf{Second pass.} Selected frames only: background removal, face landmarks and blendshape scores; 1,518 values per frame};
      \node[half, below=of extract, xshift=-2.6cm] (seq) {Ordered frame sequence\\$\downarrow$\\LSTM, with an averaging or attention readout};
      \node[half, below=of extract, xshift=2.6cm] (pool) {Mean, concatenation, mean and change, or mean and self-similarity\\$\downarrow$\\MLP, SVM, Random Forest};
      \node[box, anchor=north] (out) at ([yshift=-0.45cm]pool.south -| extract.south) {Prediction: Highly-Engaged (level 3) or not (levels 0--2)};
      \draw[arr] (clip) -- (scan);
      \draw[arr] (scan) -- (select);
      \draw[arr] (select) -- (extract);
      \draw[arr] (extract.south -| seq.north) -- (seq.north);
      \draw[arr] (extract.south -| pool.north) -- (pool.north);
      \draw[arr] (seq.south) -- (seq.south |- out.north);
      \draw[arr] (pool.south) -- (pool.south |- out.north);
   \end{tikzpicture}
   \caption{Pipeline of the proposed study. All conditions share the same clips, features, and classifiers, and differ only in the frame selection rule, the number of frames $k$, and how the selected frames are combined.}
   \label{fig:pipeline}
\end{figure}


%%
%% --- 4.2 Data --- %%
%%

\section{Data}
\label{sec:data}

\subsection{Dataset}

EngageNet \cite{singh2023} consists of about 11,300 ten-second video clips of 127 participants. Each clip is annotated with one of four engagement levels: Not-Engaged, Barely-engaged, Engaged, and Highly-Engaged, mapped to the integers 0 to 3 in that order. Clips in which the participant is absent are labeled SNP (Subject Not Present). These carry no engagement level and are removed.

The official EngageNet partition is used without re-partitioning. It is subject-independent: no participant appears in more than one split, so a classifier cannot score well by recognizing participants it saw during training. Table~\ref{tab:splits} summarizes the partition. The preliminary study (Section~\ref{sec:prelim}) drew on the full validation and test splits but only part of the training split, namely 2,925 clips from 32 of its 91 participants, because the download of the training split was incomplete. After feature extraction and the restriction described in Section~\ref{sec:commonset}, its training data contained 2,719 clips from 31 participants. The proposed study will use all three splits in full, including all 91 training participants, and will balance the training data with SMOTE (Section~\ref{sec:training}), which the preliminary study's main results did not use. Its results will therefore not be directly comparable with those of the preliminary study. Comparisons between conditions within the proposed study are unaffected, because every condition, including SME and $k$-Changepoint selection, will be re-run on the same data.

\begin{table}[t]
   \centering
   \caption[EngageNet partition]{EngageNet partition. The proposed study uses the official partition in full. Labeled clips include SNP clips, which are removed. The preliminary-study columns describe its common clip set (Section~\ref{sec:commonset}) and the participants with at least one clip in it.}
   \label{tab:splits}
   \small
   \begin{tabular}{lrrrrr}
      \toprule
      & \multicolumn{3}{c}{Official partition (proposed study)} & \multicolumn{2}{c}{Preliminary study} \\
      \cmidrule(lr){2-4} \cmidrule(lr){5-6}
      Split & Participants & Labeled clips & SNP clips & Participants & Clips used \\
      \midrule
      Training & 91$^{a}$ & 7,983 & 104 & 31 & 2,719 \\
      Validation & 11 & 1,071 & 0 & 11 & 867 \\
      Test & 26 & 2,257 & 1 & 25 & 2,005 \\
      \bottomrule
   \end{tabular}

   \vspace{0.5ex}
   {\footnotesize $^{a}$Counted from the participant numbers in the clip file names. \citeA{singh2023} report 90 training participants. The source of the difference is not known.\par}
\end{table}

\subsection{Classification Task}

The four engagement levels are reduced to a binary task with a threshold of 3. Highly-Engaged clips (level 3) form the positive class, and clips at levels 0 to 2 form the negative class. This formulation is retained from the preliminary study, where the training data used was close to balanced under it (1,375 negative and 1,344 positive clips). The full official training split is also close to balanced, with 4,139 negative and 3,740 positive clips after SNP clips are removed. Section~\ref{sec:binarization} explains why this threshold was chosen over the other two. The four-level task is outside the scope of this study.

\subsection{Common Clip Set}
\label{sec:commonset}

A clip enters the comparison only if features were extracted for it under every selection condition, and the same clips are used in every condition for training, validation, and testing. Without this restriction, a difference between two conditions could come from scoring different clips rather than from selecting different frames. In the preliminary study, the common set contained 2,719 training, 867 validation, and 2,005 test clips. A clip can be lost when none of its frames passes the validity check (Section~\ref{sec:features}) or when extraction fails. The proposed extraction will log the outcome for every clip (features extracted, no valid frame, time-out, or error), so that the common set and every exclusion from it can be reported.


%%
%% --- 4.3 Feature Extraction --- %%
%%

\section{Feature Extraction}
\label{sec:features}

Features are extracted with MediaPipe \cite{lugaresi2019}. Its Face Landmarker model returns 478 facial landmarks and 52 blendshape scores per face, and the lite version of its Pose Landmarker model returns body keypoints. Extraction runs in two passes over each clip.

\textbf{First pass.} Every second frame is read, resized so that its longer side is at most 640 pixels, and passed to both models. A frame is \emph{valid} when a face is detected and, if a body pose is also detected, two further checks hold. First, the left and right shoulders must be either both visible or both not visible, using a visibility score threshold of 0.5. Second, when both shoulders are visible, a landmark on the nose must lie horizontally between them. These checks discard frames in which the detected face and body are inconsistent with each other. For every scanned frame, eight upper-body keypoints (both shoulders, elbows, wrists, and hips) are recorded, each with $x$, $y$, and $z$ coordinates and a visibility score, giving 32 pose values. For valid frames, the two-dimensional facial landmark coordinates are also kept for the two content-based rules (Section~\ref{sec:selection}).

\textbf{Second pass.} Only the selected frames are read again, at full resolution. The background of each frame is removed with the rembg tool \cite{rembg}, using its U$^2$-Net human-segmentation model \cite{qin2020}, and the Face Landmarker is run again on the result. This gives the 52 blendshape scores and the $x$, $y$, and $z$ coordinates of the 478 landmarks (1,434 values). The 32 pose values of a selected frame are taken from the first pass.

Each selected frame is therefore described by 1,518 values: 52 blendshape scores, 1,434 landmark coordinates, and 32 pose values. Values are missing when no pose is detected in the first pass or no face is detected after background removal, and these are set to zero. Section~\ref{sec:indicators} describes what these features capture and how blendshape scores differ from action units.

Two changes will be made to the extraction used in the preliminary study. First, the outputs of the first pass (the list of valid frames, the landmark positions, and the pose values) will be stored once per clip and reused by every selection rule. All rules then choose from exactly the same valid frames, and the first pass is not repeated for each condition. Second, the outcome of every clip will be logged, as described in Section~\ref{sec:commonset}, together with the time each pass takes, from which the processing time in Section~\ref{sec:metrics} is computed.


%%
%% --- 4.4 Frame Selection Rules --- %%
%%

\section{Frame Selection Rules}
\label{sec:selection}

Let $V = (f_1, f_2, \dots, f_n)$ be the valid frames of a clip in temporal order. Each rule selects up to $k$ frames from $V$ and returns them in temporal order. When $n \le k$, every rule keeps all $n$ valid frames, so some clips contribute fewer than $k$ frames. Shorter sequences are padded and masked during training (Section~\ref{sec:training}).

\textbf{SME} keeps the first, middle, and last valid frames: $f_1$, $f_{\lfloor n/2 \rfloor + 1}$, and $f_n$. It is used only at $k = 3$.

\textbf{Positional} selection divides $V$ into $k$ segments of equal length and keeps the center frame of each, that is, $f_{\lfloor (j - 1/2)\,n/k \rfloor + 1}$ for $j = 1, \dots, k$. In words, the valid frames are split into $k$ equal parts, and the frame nearest the middle of each part is kept. This follows the segment-based sampling of Temporal Segment Networks \cite{wang2016}, and it extends position-based selection to any frame budget. Unlike SME, it does not take the first or last valid frame unless the clip has at most $2k$ valid frames, so comparing the two at $k = 3$ tests whether including the clip boundaries matters.

\textbf{Random} selection draws $k$ frames uniformly from $V$ without replacement. Because it uses no information about the clip, it estimates the accuracy that an arbitrary choice of $k$ valid frames achieves. Five draws are made, using random seeds 42 to 46. Each draw puts the valid frames of a clip in a random order once, and the rule keeps the first $k$ frames of that order. Each set of $k$ frames is still a uniform random draw, but within a draw the frames kept at $k = 3$ are also kept at $k = 5$ and $k = 7$. A larger budget therefore only adds frames, as under the two content-based rules, and fewer distinct frames must be extracted. Each draw at each $k$ produces its own feature set.

\textbf{$k$-Changepoint} selection keeps the frames at which the face moves the most. Let $\mathbf{p}_{t,i}$ be the two-dimensional position of landmark $i$ in valid frame $t$. Positions are first centered within each frame, which removes movement of the face as a whole across the image, such as a shift of the camera. The velocity between consecutive valid frames is
\begin{equation}
   v_t = \sqrt{\sum_{i=1}^{478} \left\lVert \tilde{\mathbf{p}}_{t+1,i} - \tilde{\mathbf{p}}_{t,i} \right\rVert^2},
   \qquad
   \tilde{\mathbf{p}}_{t,i} = \mathbf{p}_{t,i} - \frac{1}{478} \sum_{j=1}^{478} \mathbf{p}_{t,j}.
   \label{eq:velocity}
\end{equation}
The velocity signal is smoothed with a five-frame moving average, which replaces each value with the mean of itself and the two values on either side, so that a single jittery frame does not create a peak. Its local maxima are then found under a minimum separation of seven scanned frames (14 video frames), which prevents several frames of one continuous movement from being selected as separate peaks. The $k$ peaks with the highest velocity are kept. If fewer than $k$ peaks exist, all of them are kept, so under this rule a clip can contribute fewer than $k$ frames even when it has more than $k$ valid frames. The number of frames each clip contributes is therefore reported for every condition (Section~\ref{sec:metrics}).

Section~\ref{sec:changepoint} shows that, before smoothing, the velocity signal is equivalent to the simplest configuration of the window-sliding changepoint method reviewed by \citeA{truong2020}, with windows one frame wide. With windows that narrow, a peak marks a moment of fast movement rather than a lasting change in the signal, so the procedure is best described as velocity-peak selection.

\textbf{$k$-BOCPD} selection keeps the frames at which the face moves to a new position and stays there, as when a student leans back or looks down at notes for a while. It uses Bayesian online changepoint detection (BOCPD) \cite{adams2007}, introduced in Section~\ref{sec:changepoint}, in three steps.

\emph{Signal.} Each valid frame is described by the centered landmark positions $\tilde{\mathbf{p}}_{t,i}$ of Equation~\ref{eq:velocity}, 956 values in all. Within each clip, these are reduced to three numbers per frame: the first three principal components, which are the three directions in which the clip's landmark positions vary most. Three numbers are few enough for the method to learn their usual values from the few frames of a short segment. All three are divided by the standard deviation of the first, so that the main movement stays the largest and small jitter is not enlarged.

\emph{Score.} BOCPD treats the clip as a series of segments, stretches in which the face stays in roughly the same position. Within a segment, each of the three numbers is assumed to vary randomly around a usual value, following a Gaussian distribution whose mean and variance are learned from the segment's frames. Before a segment has any frames, a weak starting guess is used: a normal--inverse-gamma prior with $\mu_0 = 0$, $\kappa_0 = 1$, $\alpha_0 = 2$, and $\beta_0 = 1$, which a few real frames quickly override. With this prior, how well each new frame fits a segment can be computed exactly \cite{adams2007}. The \emph{hazard}, the chance that a new segment begins at any frame, is $1/20$. Segments then last 20 scanned frames on average, about 1.3 seconds when every scanned frame is valid, and a clip of 150 valid frames is expected to contain about seven changes, as many as the largest frame budget. As the frames are read, BOCPD keeps track of the \emph{run length}, the number of frames since the current segment began, as a probability for each possible value. Run lengths longer than 50 scanned frames, two and a half times the average segment, are not tracked, which keeps the method's memory short.

The score of a valid frame $f_t$ is the probability that a new segment began exactly at $f_t$. A single frame is weak evidence, because a blink or a quick glance also makes one frame look different, so the score waits for the next $L = 7$ frames:
\begin{equation}
   s^{\rightarrow}_t = P\left(r_{t+L-1} = L \mid f_1, \dots, f_{t+L-1}\right),
   \label{eq:bocpdscore}
\end{equation}
where $r_u$ is the run length at $f_u$, counting $f_u$ itself. In words, once the seven frames from $f_t$ onward have been seen, $s^{\rightarrow}_t$ is the probability that the current segment began at $f_t$. If the face moved at $f_t$ and stayed, the following frames fit a new segment and the score is high. If it moved and came back, the following frames fit the old segment again and the score is lower. The same score is also computed with the frames in reverse order, giving $s^{\leftarrow}_t$, so that $f_t$ is also judged from the seven frames before it. The two are averaged: $s_t = (s^{\rightarrow}_t + s^{\leftarrow}_t)/2$. Read in one direction only, the method is too ready to place a change near the start of a clip, where it has seen few frames. Reading in both directions, with the short memory, lets every frame be judged in the same way. A frame can therefore be scored only if it has at least $L$ valid frames on each side. $L$ equals the minimum separation used for velocity peaks.

\emph{Selection.} The frames are then chosen as for velocity peaks: the local maxima of $s_t$ are found under a minimum separation of seven scanned frames, and the $k$ highest are kept, in temporal order. The frames kept at $k = 3$ are therefore also kept at $k = 5$ and $k = 7$. A clip with more than $k$ but fewer than $2L = 14$ valid frames has no frame that can be scored, and it leaves the common clip set (Section~\ref{sec:commonset}). The prior, the hazard, the limit of 50 frames, and $L$ are fixed here, before any results are seen, and are not tuned.


%%
%% --- 4.5 Classifiers and Training --- %%
%%

\section{Classifiers and Training}
\label{sec:training}

Four classifiers are used. The LSTM receives the selected frames as an ordered sequence, and the other three receive a single vector that combines them.

\begin{itemize}
   \item The \textbf{LSTM} \cite{hochreiter1997} receives the ordered sequence of frames, padded with zeros to length $k$. It has one recurrent layer of 64 units. Its outputs are combined over the real (unpadded) time steps, by averaging or by attention (described below), then passed to a 32-unit layer with ReLU activation and a two-unit output layer.
   \item The \textbf{MLP} receives the clip's combined vector, computed over real frames only. It has layers of 64 and 32 units with ReLU activation and a two-unit output layer. Its widths match the LSTM so that it serves as a non-sequential comparison for it.
   \item The \textbf{SVM} \cite{cortes1995} uses a radial basis function kernel with the scikit-learn defaults: $C = 1$, which sets how heavily errors on the training data are penalized against keeping the gap between the classes wide, and $\gamma$ set to ``scale'', which sets from the spread of the features how far the influence of each training example reaches.
   \item The \textbf{Random Forest} \cite{breiman2001} uses 100 trees grown without depth limits.
\end{itemize}

The SVM and Random Forest receive the same combined vector as the MLP, standardized using the mean and standard deviation of the training split. The neural networks receive the unstandardized features. The LSTM and MLP are trained with the AdamW optimizer \cite{loshchilov2019}: learning rate $10^{-3}$, weight decay $10^{-4}$, batch size 64, and unweighted cross-entropy loss. Training runs for at most 20 epochs. The learning rate is halved once the validation loss has gone more than three epochs without improving, and training stops early after ten consecutive epochs without improvement. The model with the lowest validation loss is kept. Each neural configuration is trained with five seeds (42 to 46), which control weight initialization and the order of training batches. The SVM and Random Forest are deterministic given their inputs (random state 42) and are trained once per condition and frame combination.

\textbf{Combining the selected frames.} The MLP, SVM, and Random Forest are each run with four combinations of the $m \le k$ frames a clip contributes (Section~\ref{sec:temporal}). Let $\mathbf{x}_1, \dots, \mathbf{x}_m$ be the frame vectors in temporal order, each with 1,518 values.
\begin{itemize}
   \item \emph{Mean}: the average of the $m$ frame vectors (1,518 values). This is the default combination, used in the preliminary study, and it carries no temporal information.
   \item \emph{Concatenation}: the $m$ frame vectors placed one after another, with zeros in place of missing frames when $m < k$ ($1{,}518\,k$ values).
   \item \emph{Mean and change}: the mean, followed by the average absolute difference between consecutive frames, $\frac{1}{m-1}\sum_{j=1}^{m-1} \lvert \mathbf{x}_{j+1} - \mathbf{x}_j \rvert$, taken feature by feature (3,036 values). The change is set to zero when $m = 1$.
   \item \emph{Mean and self-similarity}: the mean, followed by the cosine similarity of every pair of frames, taken in temporal order ($1{,}518 + k(k-1)/2$ values, that is, 3, 10, or 21 similarities). Similarities involving missing frames are set to zero.
\end{itemize}
Differences and similarities are computed from features standardized with the mean and standard deviation of the training split. Without this step, the cosine similarity of almost any two frames would be close to 1, because the raw landmark coordinates are all positive and change little from frame to frame. All four combinations use only the selected frames. Neither the timing of the frames nor the landmark signals of the first pass, which the content-based rules use, is used, so that no rule gains information from frames it did not select. The similarities add only 3 to 21 values to the 1,518 of the mean, so in the SVM, whose decision depends on distances between whole vectors, they may carry little weight.

The LSTM is run with two ways of combining its outputs for the individual frames: averaging them over the real time steps, and attention, in which a learned score sets how much weight each time step receives. In a preliminary comparison, attention gave higher mean accuracy than averaging under all four sampling conditions, by 2.04 points over 20 runs. That comparison was made on the test split and was not tested for significance. In this study, both readouts are evaluated in every condition as part of the comparison of temporal representations (Section~\ref{sec:evaluation}), so no choice between them is made in advance.

The training data are balanced with SMOTE \cite{chawla2002} in every condition. SMOTE creates synthetic examples of the smaller class, each placed between a real example and one of its five nearest neighbors in the same class, until the two classes are the same size. It is applied to the training split only, after the common clip set is fixed, and never to the validation or test clips. For the SVM and Random Forest, it operates on the standardized combined vectors. For the LSTM and MLP, each clip's sequence of frames is flattened into one vector, together with its number of real frames, and restored afterward, so each synthetic sequence blends two real clips frame by frame. The MLP then combines each synthetic sequence in the same way as a real one. The SMOTE random state is fixed at 42. In the preliminary study, SMOTE had no material effect, because the training data were already nearly balanced. It is applied here so that every condition is trained on exactly balanced data.

No hyperparameter search is performed. All classifiers use the settings above in every condition, so that differences between conditions reflect frame selection and frame combination rather than tuning. The comparison between classifiers is therefore between default settings, not between best-tuned models.


%%
%% --- 4.6 Evaluation --- %%
%%

\section{Evaluation}
\label{sec:evaluation}

\subsection{Metrics}
\label{sec:metrics}

Each configuration is evaluated on the common test clips. The primary metric is macro F1, the unweighted mean of the F1 scores of the two classes, which counts errors on each class equally regardless of its size. Accuracy, per-class F1, and confusion matrices are also reported. The Matthews correlation coefficient (MCC) is reported as a summary of agreement between predicted and true classes \cite{chicco2020}. For binary labels, the Pearson correlation coefficient (PCC) between predictions and true labels equals the MCC, so the PCC values reported in the preliminary study are MCC values. Mean squared error (MSE) is also reported for consistency with earlier work in this project, but with labels of 0 and 1 it adds nothing. Each wrong prediction contributes 1 and each correct one 0, so MSE equals the error rate ($1 - \text{accuracy}$), and class-wise MSE equals the error rate within each class. Neither is interpreted separately.

Because the binary task pools the Not-Engaged, Barely-engaged, and Engaged levels into one class, the share of clips at each of these levels that are wrongly predicted as Highly-Engaged is also reported for every condition. This shows in particular how often Not-Engaged clips are mistaken for highly engaged ones. It is descriptive and is not tested for significance. The number of frames each clip contributes is also reported for every condition, because a clip can contribute fewer than $k$ frames when it has few valid frames or, under the content-based rules, few peaks (Section~\ref{sec:selection}). For every rule, the positions of the selected frames within the clip are also reported, which shows whether a content-based rule favors some part of the clip.

Processing time is reported for each condition: the time per clip of the second extraction pass, which grows with the number of frames selected, and the training time of each classifier. The first pass is shared by all conditions and is reported once. All times are measured on one computer, whose specifications are reported with them, and are descriptive only.

\subsection{Paired Significance Tests}

Because every condition is evaluated on the same clips, two configurations can be compared clip by clip. The exact McNemar test \cite{mcnemar1947} is used. It considers only the clips on which exactly one of the two configurations is correct, and tests whether each configuration is equally likely to be the correct one. This test is recommended for comparing two classifiers on a single test set \cite{dietterich1998}. Each comparison also reports the difference in accuracy, with a 95\% confidence interval from a paired bootstrap of 10,000 resamples of the test clips \cite{efron1993}. The bootstrap is \emph{paired} because each resample draws the same clips for both configurations, so the two are always compared on identical clips.

The comparisons are fixed in advance and grouped into seven families, 254 comparisons in all. The Holm--Bonferroni correction \cite{holm1979} is applied within each family to control the chance of a false positive across its tests. The first six families compare frame selection conditions, with the frames combined by their mean and the LSTM using its averaging readout. The seventh compares ways of combining the frames.

\begin{enumerate}
   \item \textbf{Content versus position:} each content-based rule ($k$-Changepoint and $k$-BOCPD) against positional selection, for each $k$ and classifier (24 comparisons).
   \item \textbf{Content versus chance:} each content-based rule against random selection, for each $k$ and classifier (24 comparisons).
   \item \textbf{Position versus chance:} positional against random selection, for each $k$ and classifier (12 comparisons).
   \item \textbf{Movement versus lasting change:} $k$-Changepoint against $k$-BOCPD selection, for each $k$ and classifier (12 comparisons).
   \item \textbf{Frame budget:} within each rule and classifier, $k = 3$ against $5$, $3$ against $7$, and $5$ against $7$ (48 comparisons).
   \item \textbf{Clip boundaries:} SME against positional selection at $k = 3$, for each classifier (4 comparisons).
   \item \textbf{Temporal representation:} within each of the thirteen selection conditions, each of the three temporal representations against the mean for the MLP, SVM, and Random Forest (117 comparisons), and attention against averaging for the LSTM (13 comparisons), 130 comparisons in total.
\end{enumerate}

\subsection{Repeated Seeds and Draws}
\label{sec:seeds}

A \emph{seed} is a number that fixes the random choices made when a neural network is trained, such as its starting weights and the order of its training batches. Retraining with a different seed gives a slightly different model. A \emph{draw} is one random choice of frames under random selection. Each draw selects different frames and therefore gives different features. A McNemar test compares one set of predictions with another. The neural networks have five sets of predictions, one per seed, and random selection has five, one per draw. Any comparison involving either one therefore has five versions. When both apply, seed $i$ is paired with draw $i$, and when two ways of combining the frames are compared, both use the same seeds and draws, paired by number. Each version is tested separately and corrected within its family, and the number of versions reaching significance is reported. A difference is described as reliable only if it is significant in all five versions. Accuracy is reported as the mean and standard deviation across the five versions. In the preliminary study, a single-seed test produced a significant result whose direction reversed in three of the five seeds (Section~\ref{sec:prelim}), which is why this rule is used.

\subsection{Equivalence}

A non-significant test does not show that two configurations perform the same. It may only mean that the test set was too small to detect a difference. To call two configurations \emph{equivalent}, the study needs a separate test, and for that it must first state how small a difference is too small to matter \cite{lakens2017}. This study sets that margin at 2 accuracy points, about 40 to 45 clips of a test set of 2,000 to 2,250 clips, before any results are seen. Two configurations are called equivalent only if the whole 95\% confidence interval of their accuracy difference lies between $-2$ and $+2$ points. When a comparison has five versions, equivalence is concluded only if it holds in all five, as for a reliable difference. A difference can be reliable and still lie within the margin, which happens when two configurations disagree on few clips. Each comparison therefore has one of four outcomes. It is a reliable difference if it meets the rule of Section~\ref{sec:seeds}, equivalent if its interval lies within the margin, a reliable difference smaller than the margin if both hold, and inconclusive if neither holds.

The margin of 2 points comes from the preliminary study, for three reasons. First, retraining the LSTM with different seeds changed its accuracy with a standard deviation of up to 1.3 points, so an effect of frame selection smaller than this would be lost in the variation between training runs. Second, the confidence interval of the accuracy difference between the SVM under 3-Changepoint sampling and under SME was 2.8 points wide, about $\pm 1.4$ points. Equivalence can only be shown when the whole interval fits inside the margin, so a margin much narrower than 1.5 points could rarely be confirmed on a test set of this size. Third, the significant sampling effects ranged from 2.39 to 4.79 points, so a margin of 2 points would not call differences of that size equivalent.

\subsection{Interpretation of Outcomes}

Table~\ref{tab:interpretation} states in advance how the main outcomes will be interpreted at each frame budget. Writing these rules before the experiments are run prevents the interpretation from being adjusted to fit the results. For the temporal representations, the results of the comparisons with the mean are reported condition by condition for each classifier.

Whether a temporal representation improves on the mean in general is decided by a separate rule. The rule asks in how many of the thirteen selection conditions the representation does better than the mean. Simply counting the conditions in which its comparison was significant would give a number whose reliability is unknown \cite{demsar2006}. The count is therefore made with the replicability analysis of \citeA{dror2017}, a method for comparing two methods across several datasets that states how far the count can be trusted. It has three steps.

\emph{Testing each condition.} For one representation and one classifier, each of the thirteen conditions gives a one-sided exact McNemar $p$-value. A small value means that the representation is probably better than the mean in that condition.

\emph{Counting.} The thirteen $p$-values are sorted from smallest to largest, $p_{(1)} \le p_{(2)} \le \dots \le p_{(13)}$. They are then checked in order: the smallest is multiplied by 13, the next by 12, and so on, and a value passes if the result is at most $\alpha$. The count $\hat{k}$ is the number of values that pass before the first one fails:
\begin{equation}
   \hat{k} = \max \left\{ u : (13 - v + 1)\, p_{(v)} \le \alpha \ \text{for every } v \le u \right\},
   \label{eq:replicability}
\end{equation}
with $\hat{k} = 0$ when even the smallest value fails. For example, if the six smallest values pass and the seventh is 0.0005, it is multiplied by 7, giving 0.0035, which passes when $\alpha = 0.005$. These checks are the Holm--Bonferroni correction applied to the thirteen conditions, so $\hat{k}$ is the number of conditions that remain significant after it. What the analysis adds is what the count means: with probability at least $1 - \alpha$, the representation is truly better in at least $\hat{k}$ conditions. This holds even though the conditions share test clips and are not independent, because the method makes no assumption about how they depend on each other \cite{dror2017}. Ten counts are made in all: three representations for each of the MLP, SVM, and Random Forest, and attention for the LSTM. Each therefore uses $\alpha = 0.05/10 = 0.005$, so that the chance that any of the ten is too high is at most 5\%. The same count is also made the other way round, for the conditions in which the mean does better than the representation.

\emph{Deciding.} Both counts are made in each of the five versions on its own (Section~\ref{sec:seeds}). A condition with five versions uses its $p$-value from that version, and a condition with a single version uses the same $p$-value every time. A representation is described as improving on the mean for a classifier only if two things hold in every version: it is better in at least seven conditions, more than half of the thirteen, and the mean is better in none. The counts are reported for every representation and classifier, whatever the outcome. One limit remains. Because each condition is tested on its own, a gain that is small in every condition may not be detected.

\begin{table}[t]
   \centering
   \caption[Planned interpretation of outcomes]{Planned interpretation of the frame selection outcomes at a given frame budget $k$, with the frames combined by their mean. ``Content'' stands for either content-based rule, and the two rules are interpreted separately. ``$>$'' means a reliable difference (Section~\ref{sec:seeds}), and ``$\approx$'' means equivalent within $\pm 2$ accuracy points.}
   \label{tab:interpretation}
   \begin{tabular}{L{4.3cm}L{9.2cm}}
      \toprule
      Outcome & Interpretation \\
      \midrule
      \multicolumn{2}{l}{\textit{Each content-based rule against positional and random selection}} \\
      Content $>$ random and content $>$ positional & The frames this rule selects carry information beyond both an arbitrary and an evenly spaced choice of frames. \\
      Content $\approx$ random, both $>$ positional & Evenly spaced frames are a poor choice, but this rule is not shown to add information beyond an arbitrary choice. \\
      Positional $>$ content & Covering the whole clip evenly is more useful than the frames this rule selects. \\
      Content $\approx$ random $\approx$ positional & No evidence that the selection rule matters at this frame budget. \\
      \midrule
      \multicolumn{2}{l}{\textit{The two content-based rules against each other}} \\
      $k$-BOCPD $>$ $k$-Changepoint & The beginnings of lasting changes are more informative frames than the moments of fastest movement. \\
      $k$-Changepoint $>$ $k$-BOCPD & The moments of fastest movement are more informative frames than the beginnings of lasting changes. \\
      $k$-Changepoint $\approx$ $k$-BOCPD & The two definitions of content select equally useful frames at this frame budget. \\
      \midrule
      Any other pattern & Reported as observed, without a general conclusion about the selection rules. \\
      \bottomrule
   \end{tabular}
\end{table}

\subsection{Best Configuration}
\label{sec:bestconfig}

A configuration is a combination of selection condition, frame combination, and classifier. To identify the best-performing configuration, all configurations are ranked by their macro F1 on the validation split, averaged over seeds or draws where these apply. The configuration ranked first is reported as the best, together with its results on the test split. Because the choice is made on the validation split, the test results are not inflated by the selection. They are descriptive, however. The best configuration is not tested against every other configuration, and it is specific to this dataset and task.


%%
%% --- 4.8 Ethical Considerations --- %%
%%

\section{Ethical Considerations}
\label{sec:ethics}

The study uses pre-existing data only. No new data will be collected, and the proponent will have no contact with the participants. Access to EngageNet was obtained through the dataset owners' request form and the EngageNet Dataset User Agreement, signed by the proponent with the adviser named as an additional researcher. The agreement permits academic, non-commercial, not-for-profit use only. It forbids redistributing, republishing, or disseminating the data in any form without the licensors' prior written approval, and it forbids modifying the database. It also requires that no still frame or video be used in a way that could cause a participant embarrassment or mental anguish, and that publications cite the dataset \cite{singh2023}.

The clips show the faces of identifiable people, so the study treats them as personal information. The Philippine Data Privacy Act of 2012 \cite{ra10173} applies to the processing of personal information in the Philippines, where the clips are stored and processed, and the study complies with it through the following measures.

\begin{itemize}
   \item \textbf{Storage.} The videos are stored only on the proponent's computer. They are excluded from version control and will not be shared or redistributed. The original clips are never altered. Features are computed from copies of the selected frames. Extracted landmark coordinates describe the geometry of each face and may still allow a person to be identified, so they are stored and handled under the same restrictions.
   \item \textbf{Publication.} No image of a participant will appear in this manuscript, in presentations, or in any other publication, since showing frames would republish the data. Figures use diagrams instead, as in Figure~\ref{fig:pipeline}.
   \item \textbf{Reporting.} Results are reported only in aggregate, over sets of clips. No result or engagement label is reported for an individual participant.
   \item \textbf{Intended use.} The study compares frame selection methods. Its classifiers are not intended for decisions about individual students, and their accuracy (about 70\% in the preliminary study) would not support such use. The study does not evaluate whether accuracy differs across demographic groups, which is a limitation for any future deployment.
\end{itemize}

Features were first extracted from part of the dataset in June and July 2026 for the preliminary study (Section~\ref{sec:prelim}), under the same user agreement but before any ethics review. The research ethics forms for studies using pre-existing data from human participants, which disclose this, are in Appendix~\ref{sec:appendixa}. They will be submitted for review, and no further processing of the clips, including tests of the new extraction code, will take place until the study is approved.
